\section{Stochastic Path Design and Practical Estimation}
\label{sec:method:path}

The choice of integration path significantly affects the attribution. We then introduce the \emph{Stochastic Threshold Path} (STP) as our default alternative.


Stochastic Threshold Path (STP) uses randomized head-wise activation times. For each head $j$, we sample a threshold $u_j \sim \mathcal{U}[0,1]$ and define the gate schedule:
\begin{equation}
\alpha_j(\tau; u_j) \;=\; \mathrm{clip}\!\left(\frac{\tau-u_j}{1-u_j},\,0,\,1\right).
\label{eq:alpha_path_stp}
\end{equation}
The denominator vanishes at $u_j=1$, but under $u_j\sim\mathcal{U}[0,1]$ we have $P(u_j=1)=0$, so the schedule is well defined almost surely (equivalently, implementations may draw $u_j\sim\mathcal{U}[0,1)$).
Intuitively, head $j$ stays at the zero baseline until $\tau = u_j$, then ramps linearly to $1$ by $\tau = 1$. Different heads ``turn on'' at different times, determined by their random thresholds. The STP attribution averages over $P$ independent threshold draws:
\begin{equation}
\mathrm{STP}\text{-}\mathrm{IG}_j(J)
\;=\;
\mathbb{E}_{\mathbf{u}}\!\left[ \mathrm{IG}_j[\boldsymbol{\alpha}(\cdot;\mathbf{u})] \right],
\label{eq:stp_ig_def}
\end{equation}
estimated via Monte-Carlo sampling.
DP can be viewed as the special case where all heads activate immediately and follow the same schedule. While STP is slightly more expensive due to path averaging, it better captures order variability among heads.

\paragraph{Definitions for theoretical analysis.}
To formalize the advantage of STP over the standard diagonal-path IG, we analyze which regions of the gate space each path can actually probe.

\emph{Diagonal path.} The standard integrated-gradients path is
\begin{equation}
\boldsymbol{\alpha}^{\mathrm{diag}}(\tau)=\tau \mathbf{1},
\qquad \tau\in[0,1].
\label{eq:diag_path}
\end{equation}
All coordinates move synchronously, so the path is restricted to the one-dimensional diagonal of $[0,1]^D$.

\emph{STP path.} For $\mathbf{u}=(u_1,\dots,u_D)$ with i.i.d.\ $u_j\sim\mathcal U[0,1]$, we use the same practical STP path as in \eqref{eq:alpha_path_stp}. For fixed $\mathbf{u}$, the corresponding pathwise STP attribution is
\begin{equation}
\mathrm{IG}^{\mathrm{STP}}_j(J;\mathbf{u})
=
\mathrm{IG}_j[\boldsymbol{\alpha}(\cdot;\mathbf{u})]
=
\int_0^1
\frac{\partial J(\boldsymbol{\alpha}(\tau;\mathbf{u}))}{\partial \alpha_j}
\dot{\alpha}_j(\tau;u_j)\,d\tau.
\label{eq:stp_ig_theory}
\end{equation}
Since $\alpha_j(\tau;u_j)$ is absolutely continuous, we have
\begin{equation}
\dot{\alpha}_j(\tau;u_j)=\frac{1}{1-u_j}\,\mathbf{1}\{\tau>u_j\}
\qquad \text{for a.e. }\tau\in[0,1].
\label{eq:stp_derivative}
\end{equation}

\paragraph{Theorem 1 (Diagonal-path blindness).}
Assume that $J:[0,1]^D\to\mathbb{R}$ is continuously differentiable. Then diagonal-path IG depends only on the values of $\nabla J$ along the synchronized diagonal
\[
\Delta=\{\tau\mathbf{1}:\tau\in[0,1]\}.
\]
Consequently, DP is invariant to any continuously differentiable interaction term that is invisible on this diagonal, even if that term is nonzero away from the synchronized path. Thus DP cannot distinguish objectives that differ only through such off-diagonal interactions.

\paragraph{Example.}
In two dimensions, the term
\[
H(\alpha_1,\alpha_2)=(\alpha_1-\alpha_2)^2\alpha_1
\]
vanishes on the diagonal $\alpha_1=\alpha_2=\tau$, and its partial derivatives also vanish there. It therefore captures an interaction that appears only when coordinates evolve asynchronously. Since DP enforces $\alpha_1=\cdots=\alpha_D$ throughout the path, it is completely insensitive to this type of off-diagonal structure.

\paragraph{Proof.}
Substituting the diagonal path $\boldsymbol{\alpha}^{\mathrm{diag}}(\tau)=\tau\mathbf{1}$ into \eqref{eq:ig_general} gives $d\alpha_j/d\tau=1$ for every $j$, hence
\[
\mathrm{IG}^{\mathrm{DP}}_j(J)
=
\int_0^1
\frac{\partial J(\boldsymbol{\alpha}^{\mathrm{diag}}(\tau))}{\partial \alpha_j}
\frac{d\alpha_j^{\mathrm{diag}}(\tau)}{d\tau}\,d\tau
=
\int_0^1 \frac{\partial J(\tau\mathbf{1})}{\partial \alpha_j}\,d\tau,
\]
which shows that DP depends only on $\nabla J$ restricted to $\Delta$. Now let $H:[0,1]^D\to\mathbb{R}$ be continuously differentiable and suppose that
\begin{equation}
\frac{\partial H(\tau\mathbf{1})}{\partial \alpha_j}=0
\qquad
\text{for all }\tau\in[0,1],\ j\in[D].
\label{eq:diag_blind_condition}
\end{equation}
Then $\mathrm{IG}^{\mathrm{DP}}_j(H)=0$ for every $j$, and by linearity of IG,
\begin{equation}
\mathrm{IG}^{\mathrm{DP}}(J+H)=\mathrm{IG}^{\mathrm{DP}}(J).
\label{eq:diag_blindness}
\end{equation}
This proves the claimed blindness property. \qed

\paragraph{Theorem 2 (STP explores full-dimensional asynchronous configurations).}
Assume that $J:[0,1]^D\to\mathbb{R}$ is continuously differentiable. Then, for each coordinate $j$, there exists a probability measure $\mu_j^{\mathrm{STP}}$ on $[0,1]^D$ such that every open neighborhood of every interior point $\boldsymbol{\alpha}^{\star}\in(0,1)^D$ receives positive $\mu_j^{\mathrm{STP}}$-mass. Consequently, expected STP attribution is not confined to a single synchronized trajectory; it aggregates $\partial_j J$ over a full-dimensional family of asynchronous gate configurations.

\paragraph{Proof.}
Define $\mu_j^{\mathrm{STP}}$ on $[0,1]^D$ by
\begin{equation}
\mu_j^{\mathrm{STP}}(A)
=
\mathbb{E}_{\mathbf{u}}\!\left[
\int_0^1
\mathbf{1}\{\boldsymbol{\alpha}(\tau;\mathbf{u})\in A\}
\dot{\alpha}_j(\tau;u_j)\,d\tau
\right]
\label{eq:stp_measure}
\end{equation}
for every Borel set $A\subseteq[0,1]^D$, where
\[
\dot{\alpha}_j(\tau;u_j)=\frac{1}{1-u_j}\,\mathbf{1}\{\tau>u_j\}
\]
for almost every $\tau\in[0,1]$ under the STP path \eqref{eq:alpha_path_stp}. We first show that $\mu_j^{\mathrm{STP}}$ is a probability measure. Taking $A=[0,1]^D$ in \eqref{eq:stp_measure} and using absolute continuity,
\[
\mu_j^{\mathrm{STP}}([0,1]^D)
=
\mathbb{E}_{\mathbf{u}}\!\left[\int_0^1 \dot{\alpha}_j(\tau;u_j)\,d\tau\right]
=
\mathbb{E}_{\mathbf{u}}\!\left[\alpha_j(1;u_j)-\alpha_j(0;u_j)\right]
=1,
\]
so $\mu_j^{\mathrm{STP}}$ is a probability measure. Moreover, by \eqref{eq:stp_ig_def}, \eqref{eq:stp_ig_theory}, and the definition \eqref{eq:stp_measure},
\begin{equation}
\mathrm{STP}\text{-}\mathrm{IG}_j(J)
=
\int_{[0,1]^D}
\frac{\partial J(\boldsymbol{\alpha})}{\partial \alpha_j}\,
d\mu_j^{\mathrm{STP}}(\boldsymbol{\alpha}).
\label{eq:stp_measure_repr}
\end{equation}

It remains to prove \eqref{eq:positive_mass}. Fix $\boldsymbol{\alpha}^{\star}\in(0,1)^D$ and let $O$ be any open neighborhood of $\boldsymbol{\alpha}^{\star}$. Since $\boldsymbol{\alpha}^{\star}\in(0,1)^D$, we have $\alpha_j^{\star}>0$ for every $j$. Choose
\begin{equation}
\tau^{\star}=\frac{1+\max_i \alpha_i^{\star}}{2},
\label{eq:tau_star_def}
\end{equation}
so that $\tau^{\star}\in(\max_i \alpha_i^{\star},1)$. Define
\begin{equation}
u_i^{\star}=\frac{\tau^{\star}-\alpha_i^{\star}}{1-\alpha_i^{\star}},
\qquad i=1,\dots,D.
\label{eq:u_star_def}
\end{equation}
Because $0<\alpha_i^{\star}<\tau^{\star}<1$, we have $0<u_i^{\star}<1$ for every $i$. A direct substitution into \eqref{eq:alpha_path_stp} gives
\[
\alpha_i(\tau^{\star};u_i^{\star})=\alpha_i^{\star}
\qquad \text{for all }i.
\]
Moreover, because $u_j^{\star}<\tau^{\star}$, \eqref{eq:stp_derivative} yields
\[
\dot{\alpha}_j(\tau^{\star};u_j^{\star})=\frac{1}{1-u_j^{\star}}>0.
\]
The map $(\tau,\mathbf{u})\mapsto \boldsymbol{\alpha}(\tau;\mathbf{u})$ is continuous at $(\tau^{\star},\mathbf{u}^{\star})$, so there exists a neighborhood $B$ of $(\tau^{\star},\mathbf{u}^{\star})$ and a constant $c>0$ such that $\boldsymbol{\alpha}(\tau;\mathbf{u})\in O$ and $\dot{\alpha}_j(\tau;u_j)\ge c$ throughout $B$. Since $\tau$ is integrated against Lebesgue measure and $\mathbf{u}$ has a density positive on $[0,1]^D$, the set $B$ has positive product measure. Plugging this into \eqref{eq:stp_measure} gives
\begin{equation}
\mu_j^{\mathrm{STP}}(O)>0.
\label{eq:positive_mass}
\end{equation}
This proves the claimed support property. \qed

\paragraph{Corollary (STP can capture off-diagonal interactions missed by DP).}
Let $H:[0,1]^D\to\mathbb{R}$ be a continuously differentiable interaction term that is invisible to DP but has strictly positive $j$-th partial derivative on some open region $O\subset(0,1)^D$. Then DP assigns zero attribution to $H$, whereas STP assigns strictly positive attribution weight to that region.

\paragraph{Proof.}
The diagonal-path statement follows directly from Theorem~1. To make the assumptions explicit, suppose that
\begin{equation}
\frac{\partial H(\tau\mathbf{1})}{\partial \alpha_j}=0
\qquad
\text{for all }\tau\in[0,1],\ j\in[D],
\end{equation}
and that there exist an open set $O\subset(0,1)^D$ and a constant $c>0$ such that
\begin{equation}
\frac{\partial H(\boldsymbol{\alpha})}{\partial \alpha_j}\ge c
\qquad
\text{for all }\boldsymbol{\alpha}\in O.
\label{eq:stp_detectable}
\end{equation}
Fix any $\boldsymbol{\alpha}^{\star}\in O$. Since $O\subset(0,1)^D$, we have $\boldsymbol{\alpha}^{\star}\in(0,1)^D$, so Theorem~2 applies and yields $\mu_j^{\mathrm{STP}}(O)>0$. Combining this with \eqref{eq:stp_detectable} yields
\[
\int_O \frac{\partial H(\boldsymbol{\alpha})}{\partial \alpha_j}\,
d\mu_j^{\mathrm{STP}}(\boldsymbol{\alpha})
\ge
c\,\mu_j^{\mathrm{STP}}(O)
>
0.
\]
Thus $\mathrm{IG}^{\mathrm{DP}}_j(H)=0$ while STP assigns strictly positive attribution weight to the region $O$. \qed

\paragraph{Remark (interpretation and practical implementation).}
Theorem~1 shows that standard IG along the diagonal path suffers from a structural synchronization degeneracy: it only observes gradients on the line $\alpha_1=\cdots=\alpha_D$. Theorem~2 shows that the practical STP used in this paper removes this restriction by spreading attribution mass over the interior of the gate cube. The corollary should therefore be read as a support result: STP can place nonzero attribution weight on interaction regions that are invisible to a single synchronized path, while still providing a smooth and numerically stable Monte-Carlo estimator.

\paragraph*{Numerical approximation.}
We approximate the integral in \eqref{eq:ig_general} using $K$ discrete steps. At $\tau_k = k/K$ for $k = 1, \dots, K$, we compute:
\begin{equation}
\mathrm{IG}_{j}[\boldsymbol{\alpha}] \approx \sum_{k=1}^{K} \frac{\partial J(\boldsymbol{\alpha}(\tau_k))}{\partial \alpha_{j}} \cdot \Delta \alpha_j^{(k)},
\qquad \Delta \alpha_j^{(k)} = \alpha_j(\tau_k)-\alpha_j(\tau_{k-1}).
\label{eq:ig_gate}
\end{equation}
For DP, $\Delta \alpha_j^{(k)} = 1/K$ for all $j$ and $k$, recovering the standard IG Riemann sum. For STP, $\Delta \alpha_j^{(k)}$ varies across heads and steps according to \eqref{eq:alpha_path_stp}.

\paragraph*{Post-processing and aggregation.}
To obtain non-negative importance scores, we optionally apply an element-wise post-processing function $\phi \in \{|\cdot|, \mathrm{ReLU}, \mathrm{identity}\}$. For each example, we compute
\begin{equation}
\hat{s}_{\ell,h} = \phi\big(\mathrm{STP}\text{-}\mathrm{IG}_{(\ell,h)}\big).
\end{equation}
The final head importance score $S_{\ell,h}$ is obtained by averaging $\hat{s}_{\ell,h}$ over $M$ dataset examples:
\begin{equation}
S_{\ell,h} = \frac{1}{M} \sum_{i=1}^{M} \hat{s}_{\ell,h}^{(i)}.
\end{equation}
Heads are ranked within each layer by sorting $\{S_{\ell,h}\}_{h=1}^{H}$ in descending order.

\begin{algorithm}[h]
\caption{HeadIG for DLLMs (STP with path averaging)}
\label{alg:headig}
\KwIn{Dataset $\mathcal{D}$; mask probabilities $\{p_t\}_{t=1}^{T}$; mask samples per probability $S$; IG steps $K$; path samples $P$; post-process $\phi$.}
\KwOut{Head importance scores $S_{\ell,h}$.}

\ForEach{example $(c,a)\in\mathcal{D}$}{
  \tcp{Build masking variants}
  \For{$t=1$ \KwTo $T$, $s=1$ \KwTo $S$}{
    Sample mask $m_{t,s}$ on completion with prob $p_t$; compute masked input $x^{(t,s)}$\;
  }
  \tcp{Integrated gradients with STP}
  $\mathrm{IG} \gets \mathbf{0}$\;
  \For{$p=1$ \KwTo $P$}{
    Sample thresholds $u_j \sim \mathcal{U}[0,1]$ for each head $j$\;
    $\boldsymbol{\alpha}_{\mathrm{prev}} \gets \mathbf{0}$\;
    \For{$k=1$ \KwTo $K$}{
      $\boldsymbol{\alpha}_{\mathrm{now}} \gets \boldsymbol{\alpha}(k/K; \mathbf{u})$ via \eqref{eq:alpha_path_stp}\;
      $\Delta \boldsymbol{\alpha} \gets \boldsymbol{\alpha}_{\mathrm{now}} - \boldsymbol{\alpha}_{\mathrm{prev}}$; $\boldsymbol{\alpha}_{\mathrm{prev}} \gets \boldsymbol{\alpha}_{\mathrm{now}}$\;
      Compute $\mathcal{L}$ on all variants with gated forward; $\mathbf{g} \gets \nabla_{\boldsymbol{\alpha}} (-\mathcal{L})$\;
      $\mathrm{IG} \gets \mathrm{IG} + \mathbf{g} \odot \Delta \boldsymbol{\alpha}$\;
    }
  }
  $\mathrm{IG} \gets \mathrm{IG}/P$; apply $\phi$; accumulate for averaging\;
}
\Return{$S_{\ell,h}$ by averaging over $\mathcal{D}$.}
\end{algorithm}


